% !PRINTED 38
\weberpage{38}
\begin{center}
\S\,11. \textbf{\gesperrt{Die drei Gattungen elliptischer Integrale.}}
\end{center}
Es sei jetzt
\weq{1}{f(x) = a_0x^4 + a_1x^3 + a_2x^2 + a_3x + a_4}
eine ganze Funktion dritten oder vierten Grades, und
\weq{2}{\Omega(x) = \int \frac{\Phi(x)\,dx}{\sqrt{f(x)}} = \int R(x)\,dx,}
wenn zur Abkürzung
\weq{3}{R(x) = \frac{\Phi(x)}{\sqrt{f(x)}}}
gesetzt ist, ein elliptisches Integral. Es bedeutet darin $x$ eine unbeschränkt veränderliche (auch komplexe) Größe, und $\sqrt{f(x)}$ hat für jeden Wert von $x$, für den $f(x)$ nicht verschwindet, zwei entgegengesetzte Werte. Wir verstehen unter einem \gesperrt{Punkt} nicht einen Wert von $x$ allein, sondern ein zusammengehöriges Wertepaar von $x$ und $\sqrt{f(x)}$, also einen Wert von $x$ mit einem bestimmten Vorzeichen von $\sqrt{f(x)}$. Ist daher $x_0$ irgend ein Wert von $x$, so gibt es zwei Punkte $x_0$, wenn $f(x_0)$ von Null verschieden ist, aber nur einen, wenn $f(x_0) = 0$ ist. Für $x = \infty$ bestimmen wir die Punkte nach dem Vorzeichen von $\sqrt{f(x)}:x^2$, und wir haben also zwei Punkte $x = \infty$, wenn $f(x)$ vom vierten, und nur einen, wenn $f(x)$ vom dritten Grade ist.

Nach bekannten Sätzen der Funktionentheorie gibt es, wenn $f(x_0)$ von Null verschieden ist, eine in einem gewissen Bereich konvergente Entwickelung nach dem \gesperrt{Taylorschen} Lehrsatz:
\weq{4}{R(x) = A_m(x-x_0)^m + A_{m+1}(x-x_0)^{m+1} + A_{m+2}(x-x_0)^{m+2} + \cdots.}
worin $m$ eine positive oder negative ganze Zahl oder auch Null ist, während $A_m, A_{m+1}, A_{m+2}, \dots$ Konstanten bedeuten, von denen $A_m$ nicht verschwindet. Dies ist eine Entwickelung nach steigenden Potenzen. Ergänzend tritt hinzu, falls $f(x)$ vom vierten (nicht vom dritten) Grade ist, die Entwickelung nach \gesperrt{fallenden} Potenzen
\weq{5}{R(x) = C_m x^{-m-2} + C_{m+1}x^{-m-3} + C_{m+2}x^{-m-4} + \cdots,}
von der wir sagen, daß sie in der Umgebung eines unendlich fernen Punktes gilt.
